\documentclass[11pt]{scrartcl}

\usepackage[stixtwo]{fontsetup}
\usepackage{amsmath}

\begin{document}
\begin{enumerate}
	\item A nevező gyöktényezős alakban:
	      \[
		      x^2 + x - 6 = (x + 3)(x - 2)
		      \text.
	      \]
	      A nullával való osztást nem értelmezzük, ezért $D_f = \mathbb{R}
		      \setminus \{-3, 2\}$.

	      A függvény $D_f$ minden pontjában folytonos, mert folytonos függvények
	      hányadosa.

	\item Az $x = -3$ pontban:
	      \[
		      \lim_{x \to -3} f(x)
		      = \lim_{x \to -3}
		      \underbrace{\frac{\sin\left(x^2 - 9\right)}{x^2 - 9}}_{\to 1} \cdot
		      \underbrace{\frac{x (x - 3)}{x - 2}}_{\to \frac{(-3) \cdot (-6)}{-5}}
		      = -\frac{18}{5},
	      \]
	      a bal és a jobb oldali határérték is $-\frac{18}{5}$, így itt
	      megszüntethető (elsőfajú) szakadása van a függvénynek.

	      Az $x = 2$ pontban:
	      \[
		      \lim_{x \to 2 \pm 0} f(x)
		      = \lim_{x \to 2 \pm 0}
		      \frac{
			      \overbrace{x \sin\left(x^2 - 9\right)}^{\to 2 \sin(-5) > 0}
		      }{
			      \underbrace{(x + 3)(x - 2)}_{\to 0 \pm 0}
		      }
		      = \pm \infty,
	      \]
	      vagyis a bal oldali határérték $-\infty$, a jobb oldali $+\infty$, így
	      itt másodfajú szakadása van a függvénynek.
\end{enumerate}
\end{document}
